\chapter{INTRODUÇÃO}

O impacto da pandemia no Brasil foi expressivo. Em março de 2023, o país ultrapassou a marca de 700 mil mortes causadas pela Covid-19, consolidando-se entre as nações mais afetadas pela crise sanitária mundial \cite{MS:2023}. Além das perdas humanas, a pandemia provocou forte pressão sobre a infraestrutura hospitar.

Dessa forma, a pandemia de Covid-19 fez surgir a necessidade de distribuir o material em tempo recorde e a importância de uma estratégia logística eficiente na distribuição de vacinas em cenários de emergência sanitária. Em situações pandêmicas, a rapidez na distribuição de imunizantes torna-se um fator diretamente relacionado à redução da transmissão de doenças, à mitigação do colapso dos sistemas de saúde e à preservação de vidas humanas.

No processo de imunização da população, enquanto determinadas regiões podem apresentar excesso temporário de doses armazenadas, outras podem enfrentar escassez de imunizantes. Diante dessa desigualdade de necessidade, a redistribuição eficiente de vacinas entre municípios e postos de vacinação torna-se crucial para garantir a cobertura vacinal adequada e evitar o desperdício de doses. A falta de uma logística eficaz para a redistribuição pode resultar em retenção de vacinas em locais com baixa demanda, levando ao vencimento das doses e à perda de oportunidades de imunização em regiões com alta demanda.

Há registros de grandes perdas de doses de vacinas não utilizadas. De acordo com a matéria de Craide \cite{art:CRAIDE}, entre 2021 e 2022, houve perda de 4.062.119 doses de vacinas entre 407 municípios de Minas Gerais, seguida de 3.462.098 doses desperdiçadas entre 203 municípios da Bahia, e 2.520.079 doses perdidas entre 206 municípios do Rio Grande do Sul. Além disso, de acordo com dados do Tribunal de Contas da União, no mesmo período, mais de 28 milhões de doses foram perdidas no Brasil devido ao vencimento do prazo de validade, gerando prejuízos superiores a 1 bilhão de reais. 

Segundo outra matéria \cite{art:BRASILDEFATO}, em 2022, 100 mil doses de AstraZeneca ficaram retidas, mesmo com prazo de validade próximo, e com grandes chances de serem desperdiçadas. Ainda que tenha sido mencionada uma estratégia de aumentar a faixa etária para a campanha de vacinação, não foi citada uma estratégia de redistribuição dessas doses para locais com demanda imediata.

Dessa forma, denota-se uma falta de logística para a redistribuição de vacinas, que ficam retidas até expirarem seu tempo de validade.

Nesse contexto, torna-se fundamental o desenvolvimento de ferramentas que auxiliem a tomada de decisão na redistribuição de vacinas entre localidades com diferentes níveis de oferta e demanda, especialmente em cenários de emergência sanitária. Problemas dessa natureza envolvem múltiplas restrições simultâneas, incluindo disponibilidade de doses, demandas e ofertas regionais, capacidades de transporte, limitações temporais e risco de expiração dos imunizantes. E a utilização de modelos matemáticos de otimização se justifica por sua capacidade de fornecer soluções eficientes para esse tipo de problema.

Nesse contexto, este trabalho propõe modelos de otimização baseados em Programação Linear Inteira para apoiar a distribuição de vacinas entre municípios e postos de vacinação, considerando restrições logísticas e temporais. O principal objetivo consiste em minimizar o tempo de distribuição das vacinas, buscando aumentar a eficiência logística do sistema e contribuir para respostas mais rápidas em campanhas de imunização em larga escala.

Logo, este trabalho objetivou minimizar o tempo de distribuição de vacinas e mitigar o desperdício das mesmas por meio da modelagem de um Problema de Programação Linear Inteira (PPLI) de forma a satisfazer a demanda de vacinas o máximo possível. Com esse modelo, entidades executivas, como o Munistério da Saúde e Prefeituras, poderiam mitigar os efeitos devastadores de uma pandemia, incluindo mortes, sobrecarga de infraestrutura hospitar e desperdício de vacinas.

Quanto à metodologia, inicialmente é desenvolvido um modelo intermunicipal, responsável pela redistribuição de vacinas entre municípios com excesso de oferta e municípios com demanda. Em seguida, é proposto um modelo intramunicipal, voltado à distribuição das doses entre postos de vacinação dentro de cada município. Posteriormente, ambos os modelos são integrados em formulações unificadas, permitindo representar de forma mais abrangente a cadeia logística de distribuição. As formulações consideram restrições de oferta, demanda, capacidade de transporte e viabilidade temporal associada à validade das vacinas. Os modelos são implementados computacionalmente em Python utilizando o solver Gurobi Optimizer, possibilitando a realização de experimentos computacionais e a análise dos resultados de cenários simulados para avaliação da eficiência das soluções obtidas.

Por fim, este trabalho também se relaciona ao Objetivo de Desenvolvimento Sustentável 3 (ODS 3) da Organização das Nações Unidas, voltado à promoção da saúde e do bem-estar. Ao propor mecanismos de otimização para distribuição de vacinas, busca-se contribuir para o fortalecimento de estratégias de imunização e para o aumento da eficiência logística em sistemas de saúde pública.

\bigskip
